% Structural inputs from the uploaded sources:
% Hilscher--Zeidan prove that symmetric three-term recurrence equations have a natural symplectic structure. 
% Hilscher develops oscillation theory for discrete symplectic systems with nonlinear dependence on the spectral parameter under a monotonicity condition. 
% Kratz--Hilscher prove a generalized index theorem for monotone matrix-valued functions under constant-rank hypotheses. 
% Schulz-Baldes proves a rotation-number/eigenphase theorem for finite block Jacobi matrices, including analyticity and monotonicity of eigenphases. 
% Ammann proves relative oscillation theory for finite Jacobi matrices under arbitrary perturbations, identifying differences of eigenvalue counts with weighted sign-changes of Wronskians. 

\documentclass[11pt]{article}

\usepackage[a4paper,margin=1in]{geometry}
\usepackage{amsmath,amssymb,amsthm,mathtools}
\usepackage[T1]{fontenc}
\usepackage{lmodern}
\usepackage{microtype}
\usepackage{hyperref}

\numberwithin{equation}{section}

\newtheorem{theorem}{Theorem}
\newtheorem{hypothesis}{Hypothesis}
\newtheorem{remark}{Remark}

\title{A Symplectic--Jacobi Proof Scheme for the Connectedness Threshold}
\author{}
\date{}

\begin{document}
\maketitle

Let
\[
A_n(a)=\operatorname{tridiag}(a,0,1),\qquad 0<a<1,\qquad n\ge 2,
\]
and for \(x\in\mathbb R\) set
\[
H_n(x):=(xI_n-A_n(a))^{\mathsf T}(xI_n-A_n(a)).
\]
Write
\[
\mu_1^{(n)}(x):=\lambda_{\min}(H_n(x)).
\]
Let \(\beta_n(a)\) denote the real-axis barrier already introduced in the paper, so that the
previously established barrier criterion reads
\[
\Sigma_{\varepsilon,n}\ \text{is connected}
\iff
\Sigma_{\varepsilon,n}\ \text{is simply connected}
\iff
\beta_n(a)<\varepsilon.
\]
The purpose of the present note is to record a proof scheme showing that, once two
matrix-specific algebraic lemmas are verified for the Hatano--Nelson matrix, the connectedness
threshold theorem follows from discrete symplectic oscillation theory.

\begin{hypothesis}[Folding lemma]\label{hyp:folding}
For every \(n\ge2\), every \(x\in\mathbb R\), and every \(\lambda\in\mathbb R\), the equation
\[
(H_n(x)-\lambda I_n)v=0
\]
is equivalent, after a fixed local block change of variables, to a finite symmetric three-term
recurrence
\[
S_kY_{k+1}-T_k(x,\lambda)Y_k+S_{k-1}^{\mathsf T}Y_{k-1}=0,
\qquad k=1,\dots,m_n,
\]
with \(S_k\) invertible and \(T_k(x,\lambda)\) symmetric. Equivalently, it is represented by a
finite discrete symplectic system
\[
Z_{k+1}=\mathcal S_{n,x,k}(\lambda)\,Z_k,
\qquad k=1,\dots,m_n,
\]
in the sense of Hilscher--Zeidan.
\end{hypothesis}

\begin{hypothesis}[Monotone spectral-parameter lemma]\label{hyp:monotone}
In the symplectic system of Hypothesis~\ref{hyp:folding}, the dependence on \(\lambda\) is
monotone in the sense that
\[
\Delta_{n,x,k}(\lambda)
:=J\,\dot{\mathcal S}_{n,x,k}(\lambda)\,\mathcal S_{n,x,k}(\lambda)^{-1}
\succeq 0,
\qquad k=1,\dots,m_n,
\]
and the constant-rank/image hypotheses required by the generalized index theorem are satisfied at
every singular crossing.
\end{hypothesis}

\begin{hypothesis}[One-step extension lemma]\label{hyp:extension}
For every fixed \(x\in\mathbb R\), the folded system attached to size \(n+1\) is a one-step
principal-submatrix extension of the folded system attached to size \(n\). In particular, after
folding, the pair \((n,n+1)\) fits into the finite Jacobi comparison framework to which relative
oscillation theory applies.
\end{hypothesis}

\begin{theorem}[Connectedness threshold theorem under the folding and extension lemmas]
\label{thm:conditional-threshold}
Assume Hypotheses~\ref{hyp:folding}--\ref{hyp:extension}. Then, for every
\(a\in(0,1)\) and every \(\varepsilon>0\),
\[
N_c^{\mathrm{first}}(a,\varepsilon)=N_c^{\mathrm{event}}(a,\varepsilon).
\]
More precisely,
\[
N_c^{\mathrm{first}}(a,\varepsilon)
=
N_c^{\mathrm{event}}(a,\varepsilon)
=
\min\{n\ge2:\beta_n(a)<\varepsilon\}.
\]
\end{theorem}

\begin{proof}[Proof sketch]
The previously established topological reduction already shows that, for each \(n\),
\[
\Sigma_{\varepsilon,n}\ \text{is connected}
\iff
\beta_n(a)<\varepsilon.
\]
Therefore it suffices to prove the global barrier monotonicity
\[
\beta_{n+1}(a)\le \beta_n(a),
\qquad n\ge2,\quad 0<a<1.
\]

Fix \(a\in(0,1)\) and \(x\in\mathbb R\). For each \(n\ge2\), define the counting function
\[
\mathcal N_n(x,\lambda)
:=
\#\bigl\{\text{eigenvalues of }H_n(x)\text{ in }(-\infty,\lambda)\bigr\},
\qquad \lambda\in\mathbb R.
\]
We shall show that
\[
\mathcal N_{n+1}(x,\lambda)\ge \mathcal N_n(x,\lambda)
\qquad\text{for all }x,\lambda,
\]
which immediately implies
\[
\mu_1^{(n+1)}(x)\le \mu_1^{(n)}(x)
\qquad\text{for all }x\in\mathbb R,
\]
and therefore, after taking the maxima over the real spectral gaps,
\[
\beta_{n+1}(a)\le \beta_n(a).
\]

To obtain the monotonicity of \(\mathcal N_n\), proceed in four steps.

\medskip
\noindent
\textit{Step 1: passage to a symplectic/Jacobi problem.}
By Hypothesis~\ref{hyp:folding}, the self-adjoint eigenvalue problem
\[
(H_n(x)-\lambda I_n)v=0
\]
is equivalent to a finite symmetric three-term recurrence, and hence to a finite discrete
symplectic system. This is exactly the framework in which the symplectic structure of symmetric
three-term recurrences is available; see Hilscher--Zeidan and the monograph of Ahlbrandt--
Peterson \cite{HZ2010,AhlbrandtPeterson}.

\medskip
\noindent
\textit{Step 2: monotone spectral flow in \(\lambda\).}
By Hypothesis~\ref{hyp:monotone}, the symplectic system depends monotonically on
\(\lambda\). Therefore Hilscher's oscillation theory for discrete symplectic systems with nonlinear
dependence on the spectral parameter applies \cite{Hilscher2012}. In particular, finite eigenvalues
are counted by focal points/eigenphase crossings. At singular crossings, the generalized index
theorem of Kratz--Hilscher converts the monotone matrix-valued crossing into the corresponding
jump of the Morse index \cite{KratzHilscher2013}. Consequently, the quantity \(\mathcal N_n(x,\lambda)\)
coincides with the total number of phase crossings below \(\lambda\).

\medskip
\noindent
\textit{Step 3: rotation number interpretation.}
For finite block Jacobi matrices, Schulz-Baldes constructs a unitary Prüfer matrix
\(U_N^E\) whose eigenphases are real analytic in the energy parameter and strictly increasing,
because
\[
S_N^E=\frac{1}{i}\,(U_N^E)^*\,\partial_E U_N^E
\]
is positive definite \cite{SchulzBaldes2007}. Applied to the folded system from Step~1, this
shows that the crossings counted in Step~2 are simple and ordered. Hence \(\mathcal N_n(x,\lambda)\)
is exactly the spectral counting function of \(H_n(x)\) below \(\lambda\).

\medskip
\noindent
\textit{Step 4: comparison of sizes (n) and (n+1).}
By Hypothesis~\ref{hyp:extension}, passing from size (n) to size (n+1) corresponds, in the
folded picture, to a one-step principal-submatrix extension. Relative oscillation theory for finite
Jacobi matrices under arbitrary perturbations then applies. By Ammann's theorem, the difference
of eigenvalue counts in an interval equals the number of weighted sign-changes of the Wronskian
of suitable solutions of the two Jacobi equations \cite{Ammann2014}. In the present one-step
extension situation, this yields the comparison
\[
\mathcal N_{n+1}(x,\lambda)\ge \mathcal N_n(x,\lambda)
\qquad\text{for every }x\in\mathbb R,\ \lambda\in\mathbb R.
\]

As explained above, this implies
\[
\mu_1^{(n+1)}(x)\le \mu_1^{(n)}(x)
\qquad\text{for all }x\in\mathbb R.
\]
Taking maxima over the real spectral gaps gives
\[
\beta_{n+1}(a)\le \beta_n(a),
\qquad n\ge2.
\]

Finally, since \(\beta_n(a)\to0\) as \(n\to\infty\), the set
\[
\{n\ge2:\beta_n(a)<\varepsilon\}
\]
is nonempty for every fixed \(a\in(0,1)\) and \(\varepsilon>0\). Because it is an upper tail by
the monotonicity just proved, we conclude that
\[
N_c^{\mathrm{first}}(a,\varepsilon)
=
N_c^{\mathrm{event}}(a,\varepsilon)
=
\min\{n\ge2:\beta_n(a)<\varepsilon\}.
\]
This proves the theorem.
\end{proof}

\begin{remark}
The proof above is fully standard once the two Hatano--Nelson-specific algebraic inputs are
verified:

\begin{enumerate}
\item the folding of the five-diagonal matrix \(H_n(x)-\lambda I_n\) into a symmetric three-term
recurrence / discrete symplectic system;
\item the one-step principal-submatrix extension structure under \(n\mapsto n+1\).
\end{enumerate}

All remaining steps are supplied by the discrete symplectic oscillation literature and the relative
oscillation theory for finite Jacobi matrices.
\end{remark}

\begin{thebibliography}{99}

\bibitem{AhlbrandtPeterson}
C.~D.~Ahlbrandt and A.~C.~Peterson,
\emph{Discrete Hamiltonian Systems: Difference Equations, Continued Fractions, and Riccati Equations},
Kluwer, Dordrecht, 1996.

\bibitem{Ammann2014}
K.~Ammann,
\emph{Relative oscillation theory for Jacobi matrices extended},
Operators and Matrices \textbf{8} (2014), 99--115.

\bibitem{Hilscher2012}
R.~\v{S}.~Hilscher,
\emph{Oscillation theorems for discrete symplectic systems with nonlinear dependence in spectral parameter},
Linear Algebra Appl. \textbf{437} (2012), 2922--2960.

\bibitem{HZ2010}
R.~\v{S}.~Hilscher and V.~Zeidan,
\emph{Symmetric three-term recurrence equations and their symplectic structure},
Adv. Difference Equ. 2010, Article ID 626942.

\bibitem{KratzHilscher2013}
W.~Kratz and R.~\v{S}.~Hilscher,
\emph{A generalized index theorem for monotone matrix-valued functions with applications to discrete oscillation theory},
SIAM J. Matrix Anal. Appl. \textbf{34} (2013), 228--243.

\bibitem{SchulzBaldes2007}
H.~Schulz-Baldes,
\emph{Rotation numbers for Jacobi matrices with matrix entries},
Math. Phys. Electron. J. \textbf{13} (2007), Paper 5.

\end{thebibliography}

\end{document}
